\chapter{The Station, and what has been proved}\label{ch:station}

An ATE \textbf{Station} is one installation of ATE on one computer. As the administrator you are
responsible for three things on that computer, and they are separate:

\begin{center}
\footnotesize
\begin{tabular}{@{}>{\raggedright\arraybackslash}p{3.2cm}>{\raggedright\arraybackslash}p{5.2cm}>{\raggedright\arraybackslash}p{5.6cm}@{}}
\toprule
Part & Windows & Linux \\
\midrule
The program & \texttt{C:\textbackslash Program Files\textbackslash ATE} (one directory,
  \texttt{ate.exe} plus \texttt{\_internal}) & \texttt{/opt/ate} for the \texttt{.deb}, with
  \texttt{/usr/bin/ate} linking to it; the tarball goes where \texttt{install.sh} is told to put it \\
The config & \texttt{\%PROGRAMDATA\%\textbackslash ATE\textbackslash config.json} & \texttt{/etc/ate/config.json} \\
The data & \texttt{\%PROGRAMDATA\%\textbackslash ATE} & \texttt{/var/lib/ate} \\
\bottomrule
\end{tabular}
\end{center}

On Windows the config and the data share one directory; on Linux they are deliberately apart,
because the config is meant to be admin-only and the data has to be writable by whoever runs
Test Runs. The data directory holds the database (\texttt{ate.db}), \texttt{reports},
\texttt{logs} and \texttt{definitions}. Nothing else on the machine matters to ATE.

The words used throughout this manual are the project's: an \textbf{Operator} is the person
running a Test Run (not a Windows user account); a \textbf{UUT} is the unit under test, named by
its UUT Type and serial; a \textbf{Media Release} is a set of Lua files delivered together,
described by one \textbf{Test Definition}; a \textbf{Device} is a connection a Test Step talks
through, and \textbf{Equipment} is a record of what was on the bench, declared by the schedule so
it appears on the Test Report. They are not the same thing and ATE never links them.

\section{What has been checked, and what has not}

ATE is a desktop application that has only ever been exercised headless. Say this to anyone who
asks how much of it is proven, because the difference matters when a Station misbehaves.

\textbf{Proved.} The whole test suite runs on Windows and on Linux, on CPython 3.13 and 3.14,
with \texttt{QT\_QPA\_PLATFORM=offscreen}. Config validation, database migration and its refusals,
the permissions treatment \texttt{ate setup} applies, checksums, Device error handling, the Test
Report and its PDF helper, log export --- each is pinned by tests. The Windows installer has been
compiled with Inno Setup and installed on a clean CI runner: silent install, machine PATH entry,
the Station setup it runs, a headless scan of the example Media Release, an upgrade over the top
of the first install, and an uninstall that leaves the PATH as it found it. The Linux artifacts
have been installed in clean \texttt{debian:stable} containers --- the \texttt{.deb} with
\texttt{apt-get}, the tarball with \texttt{install.sh}, the wheel with \texttt{pipx} --- and the
built bundle has run a Test Schedule end to end there, writing a Test Report PDF that
\texttt{pdftotext} reads.

\textbf{Not proved.} A real desktop session. No window of ATE has ever been looked at by a
human: every run mentioned above had no display, so the layout, the fonts, and what a click
does are not verified by anything. Real instruments have not been connected: every Device test
uses a loopback peer, a \texttt{socket://} URL that pyserial accepts in place of a port, or
pyvisa-py's own backends. No GPIB, USBTMC or RS-232 hardware has been attached, and NI-VISA is
never shipped. The installer is unsigned, because there is no certificate yet; what a given
Windows machine will show for an unsigned installer is not established by anything in this
repository. Finally, the stall one Test Run hit on the Windows CI runner (packet 32b) is
\emph{believed} to be a slow disk, not measured: the fix is proven, the cause is not.

\section{What you will do, in order}

\begin{enumerate}
  \item Install the program (chapter~\ref{ch:installing}).
  \item Prepare the Station with \texttt{ate setup} (chapter~\ref{ch:setup}).
  \item Set the config keys that describe this bench (chapter~\ref{ch:config}).
  \item Create the database, and keep a copy of it (chapter~\ref{ch:database}).
  \item Put Media Releases where ATE can find them (chapter~\ref{ch:media-releases}).
  \item Run the station day to day (chapter~\ref{ch:running}), and know where the logs are
    (chapter~\ref{ch:logs}).
  \item When something is wrong, start from the message the machine printed
    (chapter~\ref{ch:troubleshooting}).
\end{enumerate}

\cando{say what lives where on a Station, and state honestly which of ATE's behaviours have been
tested and which have not.}
